\section{Chosen for All}

Every address has an addressee. Revelation might therefore begin anywhere, but it cannot begin nowhere. The biblical claim is more unsettling: it begins with choices. God calls Abram from his country, binds himself by covenant, frees one enslaved people, gives them a law, and causes his name to dwell among them. These are not illustrations subsequently attached to a universal religious theory. They are the history in which God makes himself known.

The universal purpose, however, is present in the first narrowing. The call of Abram contains the promise, ``In you all the families of the earth shall be blessed.'' Israel is summoned to be ``a kingdom of priests and a holy nation.'' The prophetic vocation widens into ``a light to the nations,'' so that God's salvation may reach ``to the end of the earth.'' Election is particular, but its inner direction is outward.\endnote{See Gen 12:1--3; Exod 19:3--6; Isa 49:1--6; and CCC 59--64 on Abraham, Israel, and the prophetic promise of a salvation including all nations; \sourceurl{https://www.vatican.va/content/catechism/en/part_one/section_one/chapter_two/article_1/ii_the_stages_of_revelation.html}{official English}. Scriptural phrases are brief working-English renderings; Scripture retains its independent rights status.}

That movement prevents election from becoming a theory of divine favoritism. Deuteronomy does not tell Israel that it was chosen for being larger or more impressive. It says, ``It was because the Lord loved you.''\endnote{Deut 7:6--8 places election in God's love and fidelity rather than Israel's numerical greatness. The article draws the vocation-not-boast conclusion as source-grounded synthesis rather than as a quotation from the passage.} The gift creates a vocation, not a warrant to boast. A priestly people is set apart not so that blessing may stop there, but so that it may be borne. The prophets can therefore judge the elected people by the very covenant that distinguishes them. Scripture's unsparing record of Israel's sins makes little sense as national self-congratulation.

It would be equally false to make Israel a temporary delivery system, honored only until Christianity can discard it. Jesus does not emerge from a generic humanity. He is ``born of a woman, born under the law''; the woman is a daughter of Israel, and the law is Israel's covenantal law. When Jesus says, ``salvation is from the Jews,'' the historical grammar of salvation is not an embarrassment he intends later universality to erase.\endnote{Gal 4:4--5; John 4:19--24; and CCC 422--423. The claim concerns the permanent historical grammar of the Incarnation, not ethnic ownership of salvation.}

The Church accordingly discovers her bond with the Jewish people when she examines her own mystery. To Israel belong the covenants, worship, promises, and patriarchs; from Israel according to the flesh comes the Christ. Paul will not allow Gentile Christians to turn their reception into contempt: ``the gifts and the call of God are irrevocable.'' The mystery of Israel's relation to Christ is not solved by pretending that the root has ceased to matter.\endnote{See Rom 9:4--5; 11:16--29; CCC 839--840, \sourceurl{https://www.vatican.va/content/catechism/en/part_one/section_two/chapter_three/article_9/paragraph_3_the_church_is_one\%2C_holy\%2C_catholic\%2C_and_apostolic.html}{official English}; and \emph{Nostra aetate}, no.~4. The article affirms Christian fulfillment in Christ without treating the Jewish people as a rejected or merely archaeological remainder. It does not attempt to settle every disputed account of covenant, mission, and the present relation of Israel and the Church.}

Particular election and universal purpose are thus not two stages between which God changes his mind. The one is the chosen means of the other. This does not answer every question about why this people, this time, or this road was chosen. Gratuitous love is not a theorem from which all its concrete gifts can be deduced in advance. It does show why the charge of mere provincialism is too quick. A road is not opposed to a destination because it begins somewhere.

The same pattern will reach its limit in a single Israelite. He will not offer one more local account of the divine beside many others. The Christian claim is that the one through whom all things were made has entered the history prepared for him.
